% 本节仅保留功能讲述框架。
\section{感知通信与视觉推理}
\KeyPoint{\StoryFiveTitle}{\StoryFiveQuestion}

\subsection{功能边界与输入输出}
\noindent\begin{tabularx}{\linewidth}{@{}p{0.24\linewidth}X@{}}
\toprule
功能 & 配置或状态 \\
\midrule
数据采集 & \AcquisitionFunction \\
通信 & \CommunicationFunction \\
移动数据中心 & \MobileDataCenterFunction \\
视觉处理 & \VisionFunction \\
智能推理 & \InferenceFunction \\
模型 & \AIModel \\
\bottomrule
\end{tabularx}

\subsection{任务链路与数据组织}
\FigureSlot{感知与推理任务流程图}{应串联数据来源、处理步骤、通信或存储、推理结果及使用方式，标清模块输入输出。模型、算法和执行位置不提前填写。}

\ContentSlot{功能说明与异常处理}{应逐项说明功能目的、输入条件、输出形式、任务关系与验证场景；同时说明数据缺失、通信异常或处理失败时需要考虑的行为。具体策略待规划审阅。}
